
\section{Conclusion}
In this work, we have explored over-tokenized transformer for language modeling. By systematically analyzing the impact of token granularity and vocabulary size across models of varying scales, we uncover an important scaling phenomena that inspires tokenizer design. Our findings reveal that larger input vocabularies consistently enhance model performance, no matter model scales, while larger output vocabulary can be harmful and difficult to learn for smaller models. These insights have significant implications for the continued scaling of LLMs. With such inspirations, we introduce over-encoding technique that scales up the input vocabulary via multi-gram token embeddings. Moreover, we develop over-tokenized transformer combining over-encoding and multi-token prediction. Extensive experiments validated the effectiveness of our methods, showcasing significant gains in model performance. Notably, we demonstrated a strong log-linear relationship between input vocabulary size and training loss, highlighting a new dimension in the scaling laws of LLMs. These results emphasize the importance of tokenizer design as a key driver of advancements in language modeling.

In summary, this work bridges the gap between tokenizer design and model scaling, positioning tokenization as a critical factor in the development of next-generation language models. We believe that the proposed Over-Tokenized Transformers framework and the insights derived from our study will inspire further research into tokenization strategies and their role in scaling LLMs efficiently and effectively.


\section*{Acknowledgements}

This research was conducted at ByteDance Inc. We are grateful for the suggestions and discussions from Jundong Zhou, Zihao Huang, Yu Bao, Yike Yuan, Sijun Zhang and other colleagues at ByteDance Seed.

\section*{Impact Statement}

This paper presents work whose goal is to advance the field of 
Machine Learning. There are many potential societal consequences 
of our work, none which we feel must be specifically highlighted here.
